\section{Початкова матриця простежуваності вимог}\label{sec:dodatok-matrytsia}

Цей додаток містить початкову матрицю простежуваності (Traceability Matrix), яка забезпечує двосторонній зв'язок між зафіксованими пунктами погодженого консенсусу $K^*$, формальними ідентифікаторами вимог (\texttt{REQ-F}, \texttt{REQ-TDD}, \texttt{REQ-LOG}, \texttt{REQ-NF}), конкретними умовами їх верифікації та відповідними BDD-сценаріями перевірки поведінки (\tabref{tab:traceability}).

\begin{small}
\setlength{\tabcolsep}{3.5pt}
\begin{longtable}{@{}p{2.2cm} p{1.4cm} p{1.6cm} p{7.4cm} p{2.4cm}@{}}
\caption{Матриця простежуваності: Вимоги $\rightarrow$ База $K^*$ $\rightarrow$ Критерії $\rightarrow$ BDD-сценарії}\label{tab:traceability}\\
\toprule
\textbf{ID вимоги} & \textbf{Пункт $K^*$} & \textbf{Розділ SRS} & \textbf{Критерій прийняття} & \textbf{BDD-сценарій} \\
\midrule
\endfirsthead
\multicolumn{5}{l}{\small\itshape Продовження таблиці~\ref{tab:traceability}}\\
\toprule
\textbf{ID вимоги} & \textbf{Пункт $K^*$} & \textbf{Розділ SRS} & \textbf{Критерій прийняття} & \textbf{BDD-сценарій} \\
\midrule
\endhead
\bottomrule
\endfoot
\bottomrule
\endlastfoot
REQ-F-001 & K2, K22 & \S\ref{sec:funktsional-katalog} & Публічний перегляд фільмів і сеансів без автентифікації & --- \\
REQ-F-002 & K2, K31 & \S\ref{sec:funktsional-katalog} & Відображення схеми залу; вибір місця не резервує його & --- \\
REQ-F-003 & K23 & \S\ref{sec:funktsional-auth} & Бронювання та скасування суто після автентифікації & --- \\
REQ-F-004 & K24 & \S\ref{sec:funktsional-auth} & Унікальний числовий внутрішній ID користувача & --- \\
REQ-F-005 & K24 & \S\ref{sec:funktsional-auth} & Доступ виключно до власних бронювань (захист IDOR) & --- \\
REQ-F-006 & K26 & \S\ref{sec:funktsional-auth} & Відокремлений адміністративний розділ за моделлю RBAC & --- \\
REQ-F-007 & K4 & \S\ref{sec:funktsional-bronuvannia} & Обмеження кількості місць: від 1 до 6 включно & --- \\
REQ-F-008 & K5 & \S\ref{sec:funktsional-bronuvannia} & Атомарне створення в момент підтвердження без hold & Сценарій 1 \\
REQ-F-009 & K29 & \S\ref{sec:funktsional-bronuvannia} & Рівно одне успішне бронювання при конкуренції за місце & Сценарій 1 \\
REQ-F-010 & K30 & \S\ref{sec:funktsional-bronuvannia} & Відхилення бронювання, якщо хоч одне місце зайняте & Сценарій 2 \\
REQ-F-011 & K36 & \S\ref{sec:funktsional-skasuvannia-user} & Скасування лише всього бронювання повністю (без часткового) & Сценарій 3 \\
REQ-F-012 & K37 & \S\ref{sec:funktsional-skasuvannia-user} & Скасування дозволене виключно до настання часу сеансу & Сценарії 3, 4, 5 \\
REQ-F-013 & K38 & \S\ref{sec:funktsional-skasuvannia-user} & Перехід у статус «Скасоване» та негайне звільнення місць & Сценарій 3 \\
REQ-F-014 & K39 & \S\ref{sec:funktsional-skasuvannia-user} & Ідемпотентність повторних спроб скасування & --- \\
REQ-F-015 & K3 & \S\ref{sec:funktsional-admin} & Управління довідниками фільмів, залів і сеансів & --- \\
REQ-F-016 & K44 & \S\ref{sec:funktsional-admin} & Вільне редагування параметрів сеансу без бронювань & --- \\
REQ-F-017 & K45 & \S\ref{sec:funktsional-admin} & Блокування прямого редагування сеансу з бронюваннями & Сценарій 6 \\
REQ-F-018 & K46 & \S\ref{sec:funktsional-admin} & Явне скасування сеансу адміном з обов'язковою причиною & Сценарій 6 \\
REQ-F-019 & K47 & \S\ref{sec:funktsional-admin} & Каскадний перехід бронювань у «Скасовано» та збереження & Сценарій 6 \\
REQ-F-020 & K48 & \S\ref{sec:funktsional-admin} & Відображення причини скасування сеансу в кабінеті & Сценарій 6 \\
REQ-TDD-001 & K7, K11 & \S\ref{sec:tdd-process} & Обов'язковість попереднього створення автотестів & --- \\
REQ-TDD-002 & K12 & \S\ref{sec:tdd-process} & Дотримання 5-крокового циклу Red-Green-Refactor & --- \\
REQ-TDD-003 & K13 & \S\ref{sec:tdd-scope} & Покриття тестами всіх погоджених граничних кейсів & Сценарії 1--6 \\
REQ-TDD-004 & K14 & \S\ref{sec:tdd-scope} & Якісна перевірка правил без штучного порогу покриття & --- \\
REQ-TDD-005 & K15, K43 & \S\ref{sec:tdd-scope} & Вимоги до тестового раннера (швидкість, емуляція) & --- \\
REQ-LOG-001 & K8, K16 & \S\ref{sec:loguvannia-podii} & Автоматичне логування 5 обов'язкових категорій подій & --- \\
REQ-LOG-002 & K17 & \S\ref{sec:loguvannia-podii} & Наявність обов'язкових полів аудиту в кожному записі & --- \\
REQ-LOG-003 & K18 & \S\ref{sec:loguvannia-bezpeka} & Сувора заборона збереження паролів і токенів у логах & --- \\
REQ-LOG-004 & K19 & \S\ref{sec:loguvannia-zberezhennia} & Персистентність та збереження логів після рестарту & --- \\
REQ-LOG-005 & K20 & \S\ref{sec:loguvannia-testy} & Збереження окремих звітів кожного прогону тестів TDD & --- \\
REQ-NF-001 & K5, K30 & \S\ref{sec:nefunktsional-atomarnist} & Транзакційна атомарність бронювання (ACID) & Сценарії 1, 2 \\
REQ-NF-002 & K40, K41 & \S\ref{sec:nefunktsional-atomarnist} & Персистентність та збереження узгодженості даних & --- \\
REQ-NF-003 & K40 & \S\ref{sec:nefunktsional-bezpeka} & Безумовний серверний контроль прав доступу RBAC & --- \\
REQ-NF-004 & K24 & \S\ref{sec:nefunktsional-bezpeka} & Ізоляція даних бронювань між користувачами (IDOR) & --- \\
REQ-NF-005 & K27 & \S\ref{sec:nefunktsional-bezpeka} & Заборона самопризначення ролі адміністратора клієнтом & --- \\
REQ-NF-006 & K29, K43 & \S\ref{sec:nefunktsional-konkurentnist} & Коректна конкурентна обробка без подвійних бронювань & Сценарій 1 \\
\bottomrule
\end{longtable}
\end{small}
